% !TeX root = ../main.tex

\chapter{相关工作}
\label{cha:related}

\section{自动驾驶场景理解与问答基准}
\label{sec:related-benchmark}

自动驾驶多模态研究的重要基础，是构建能够真实反映道路语义复杂性、目标距离变化和任务条件约束的评测基准。早期研究主要围绕检测、跟踪、三维重建与规划展开，数据集强调几何精度和时序一致性；随着大模型和多模态模型的发展，研究重点逐渐转向场景问答、文本读取、异常事件解释和任务导向语义推理\cite{caesar2020nuscenes,dosovitskiy2017carla,marcu2024lingoqa,tom2023readingbetweenthelanes,sima2024drivelm,ding2024holisticautonomousdriving,charoenpitaks2025tbbench,guo2024drivemllm,choudhary2023talk2bev}。这一路径推动自动驾驶研究从“能否检测目标”扩展为“能否理解场景并回答任务问题”。

从评测特征看，现有自动驾驶场景理解基准大致可以分为三类。第一类是以感知和几何建模为核心的数据集，它们为后续多模态研究提供了必要的视觉基础，但对开放语义表达支持有限。第二类是以视觉问答或场景描述为核心的数据集，它们增强了语言接口，却未必显式控制距离、目标尺度和风险等级。第三类是围绕交通任务约束构建的距离敏感问答基准，其价值在于能够把模型错误与安全相关场景条件对应起来。对本文而言，真正关键的不是“是否有问答标签”，而是“基准是否能够揭示远距离和长尾失效模式”。

与本文最相关的是 DTPQA\cite{theodoridis2025dtpqa,theodoridis2025evaluating}。该基准显式标注目标距离，并将问题限定为交通感知型问答，使研究者能够观察模型在近距离、中距离和远距离条件下的性能变化。与通用多模态问答数据集相比，DTPQA 更适合分析自动驾驶场景中的距离敏感退化，也更适合研究系统如何围绕高风险失效样本组织补救流程。因此，本文选择该基准作为主定量实验入口。

\section{自动驾驶中的多模态大模型}
\label{sec:related-advlm}

随着视觉语言模型与多模态大模型的发展，自动驾驶研究开始尝试将更强的语义理解能力引入场景感知与驾驶辅助。DriveLM、SENNA、SimpleLLM4AD、LaVIDA Drive、OmniDrive、V3LMA、LightEMMA 和 DynRSL-VLM 等工作分别从图结构问答、多传感器理解、语言辅助决策和场景级推理等角度探索了多模态大模型在驾驶任务中的潜力\cite{sima2024drivelm,jiang2024senna,zheng2024simplellm4ad,jiao2024lavidadrive,wang2024omnidrive,lubberstedt2025v3lma,qiao2025lightemma,zhou2025dynrslvlm}。这些工作普遍表明，当模型具有更强的视觉语义融合能力时，场景解释、关系建模和开放问题回答会获得显著提升。

从方法路线看，相关研究大致可归纳为两类。第一类尝试构建统一的驾驶多模态大模型，通过更大规模的多模态预训练和更复杂的输入组织来提升综合能力；第二类则将语言能力作为上层接口，用于辅助感知结果解释、场景摘要生成或策略推理。前者强调能力上限，后者强调任务接口扩展，但两者大多默认充足算力、稳定带宽和集中式推理环境，因此在部署层面往往并未真正解决边缘实时性和云端代价控制问题。

这也构成了本文与现有工作的一个重要差异。本文并不试图进一步扩大统一模型，也不将云端大模型视为默认推理主链路，而是关注在角色分层和推理路径分层条件下，如何利用边缘模型与云端模型的能力互补性获得系统收益。换言之，本文的问题不是“如何训练更强的驾驶大模型”，而是“如何在现有模型冻结的条件下组织更有效的部署链路”。

\section{视觉语言模型的感知短板}
\label{sec:related-visualgap}

大量研究表明，视觉语言模型虽然在开放语义任务上表现突出，但在基础视觉感知方面仍存在系统性脆弱性。已有工作从视觉细节识别、几何关系理解、局部区域证据利用、对象计数稳定性与视觉幻觉等角度揭示了这类模型的局限\cite{rahmanzadehgervi2024vlmsblind,tong2024eyeswideshut,gou2024imagedetails,kamoi2024visonlyqa,kaduri2025whatsintheimage,guan2024hallusionbench}。更广泛的多模态评测工作也显示，不同模型在感知能力与语言推理能力之间往往存在明显失衡\cite{liu2024mmbench,yue2024mmmu,duan2024vlmevalkit}。

这类短板对自动驾驶任务具有放大效应。首先，远距离目标和小尺度目标本就具有低分辨率、低对比度和易遮挡等视觉特征，视觉语言模型若缺少精细局部证据建模，容易产生漏检或保守回答。其次，驾驶问答常要求模型在对象存在性判断与上下文推理之间保持一致，一旦底层感知存在偏差，语言层面的解释往往会被错误视觉锚点放大。再次，真实道路场景还具有天气变化、光照变化和背景干扰等复杂因素，这进一步增加了视觉证据利用的难度。

因此，本文并不将云端协同理解为对边缘答案进行简单的语言润色，而是将其理解为对关键视觉证据的定向补充。只有当系统能够重新围绕原始图像组织观察过程时，云端协同才可能真正缓解远距离场景中的结构性失效。

\section{测试时自适应、Agent 与外部知识}
\label{sec:related-agent}

在自然语言处理和通用智能体研究中，ReAct、Reflexion、Self-Refine、Toolformer、RAG 与 REALM 等方法表明，模型能力并不一定只能通过参数更新获得增强，还可以通过外部工具、检索机制、自我反思和外部记忆在推理时得到扩展\cite{yao2023react,shinn2023reflexion,madaan2023selfrefine,schick2023toolformer,lewis2020rag,guu2020realm,wei2022cot}。这一类方法的共同思想，是将部分能力从参数空间迁移到上下文组织、知识检索和过程控制中，从而使模型在冻结权重条件下依然具备一定的适配能力。

但是，自动驾驶场景中的测试时自适应与通用文本任务存在显著差异。其一，自动驾驶系统需要输出可追溯、可审计的结构化结果，而不是仅仅生成更长的自然语言解释。其二，自动驾驶系统工作在严格的时延预算和安全约束下，外部知识调用和多轮反思都必须受到清晰的触发条件与回退规则约束。其三，自动驾驶中的错误传播代价更高，一旦不恰当的外部知识被注入当前样本，就可能造成比原始错误更严重的系统风险。

正因如此，本文并不直接移植通用智能体中的反思流程，而是将测试时自适应进一步具体化为“结构化技能沉淀 + 约束化再感知”机制。与自由文本记忆相比，这一路径更强调知识表示的边界、复用条件与安全回退策略，更符合自动驾驶系统的部署需求。

\section{边缘云协同与部署约束下的多模态系统}
\label{sec:related-system}

边缘云协同已经成为智能系统研究中的重要方向，其核心问题是如何在边缘实时性与云端高容量之间取得平衡\cite{xie2025vlmsreadyad,chen2025automatedevaluationcornercases}。在多模态基础模型方面，CLIP、Flamingo、LLaVA、InternVL、Qwen-VL、DeepSeek-VL 和 DINOv2 等工作提供了丰富的视觉表征与跨模态推理能力\cite{radford2021clip,alayrac2022flamingo,liu2023visualinstructiontuning,chen2024internvl,zhu2025internvl3,bai2023qwenvl,wang2024qwen2vl,bai2025qwen25vl,lu2024ovis,lu2024deepseekvl,wu2024deepseekvl2,oquab2024dinov2,tong2024cambrian1}；而 ViT、RoFormer、MoE、FlashAttention-2 和 Transformers 等研究则为这类系统提供了视觉编码、稀疏计算与推理效率基础\cite{dosovitskiy2021vit,su2024roformer,shazeer2017moe,dao2024flashattention2,wolf2020transformers}。

尽管如此，现有边缘云协同研究更多关注任务卸载、算力调度或模型压缩，对“何时应触发云端纠错”“如何让云端结果沉淀为后续可复用知识”讨论相对有限。换言之，现有工作通常把边缘云协同理解为算力资源的再分配问题，而本文更关注其在长尾场景中的错误修复价值与长期经验积累作用。

本文与这类工作之间的关系在于：本文继承其关于角色分层和算力分层的基本思想，但进一步把研究重点放在选择性协同触发、结构化知识外化以及证据边界清晰的实验叙事上。

\section{本文与现有工作的区别}
\label{sec:related-difference}

综合以上相关研究可以发现，现有工作已经分别从基准构建、多模态大模型、视觉短板、测试时增强和边缘云协同等不同角度讨论了自动驾驶场景理解问题，但仍缺少一个同时强调部署约束、选择性云协同和结构化知识复用的统一框架。本文的区别主要体现在以下三点。

\begin{enumerate}
  \item 在问题定位上，本文关注的是冻结主模型权重条件下的部署导向感知增强，而非通过额外训练追求离线性能极值。
  \item 在方法组织上，本文将边缘模型、云端模型与结构化技能库统一纳入同一闭环，强调角色分层、知识外化和过程可追溯。
  \item 在实验叙事上，本文将在线系统收益、结构化技能收益与真实数据外部效度严格分层，避免将不同证据强度的结论混写为统一主张。
\end{enumerate}

\section{本章小结}
\label{sec:related-summary}

综上所述，现有研究已经从基准构建、多模态大模型、视觉短板、测试时增强和边缘云协同等不同角度讨论了自动驾驶场景理解问题，但仍缺少一个同时强调部署约束、选择性云协同和结构化知识复用的统一框架。本文的研究定位正是在这一交叉点上展开：在不宣称完整基准最优的前提下，给出一种证据边界清晰、面向实际部署约束的免训练系统方法。
